  \documentclass[11pt, reqno]{amsart}
\usepackage{amsmath, amsthm, amscd, amsfonts, amssymb, latexsym}
\renewcommand{\baselinestretch}{1.5}
\parskip .5ex
\def\Xint#1{\mathchoice
{\XXint\displaystyle\textstyle{#1}}%
{\XXint\textstyle\scriptstyle{#1}}%
{\XXint\scriptstyle\scriptscriptstyle{#1}}%
{\XXint\scriptscriptstyle\scriptscriptstyle{#1}}%
\!\int}
\def\XXint#1#2#3{{\setbox0=\hbox{$#1{#2#3}{\int}$ }
\vcenter{\hbox{$#2#3$ }}\kern-.6\wd0}}
\def\ddashint{\Xint=}
\def\dashint{\Xint-}
\newcommand{\aint}{\dashint}
\newcommand{\bm}[1]{\mb{\boldmath ${#1}$}}
\newcommand{\beqa}{\begin{eqnarray*}}
\newcommand{\eeqa}{\end{eqnarray*}}
\newcommand{\beqn}{\begin{eqnarray}}
\newcommand{\eeqn}{\end{eqnarray}}
\newcommand{\bu}{\bigcup}
\newcommand{\bi}{\bigcap}
\newcommand{\iy}{\infty}
\newcommand{\bs}{\boldsymbol}
\newcommand{\lt}{\left}
\newcommand{\rt}{\right}
\newcommand{\ra}{\rightarrow}
\newcommand{\Ra}{\Rightarrow}
\newcommand{\Lra}{\Leftrightarrow}
\newcommand{\lgra}{\longrightarrow} 
\newcommand{\Lgra}{\Longrightarrow}
\newcommand{\lglra}{\longleftrightarrow}
\newcommand{\Lglra}{\Longleftrightarrow}
\newcommand{\bP}{\mathbb P}
\newcommand{\bQ}{\mathbb Q}
\newcommand{\A}{\mathbb A}
\newcommand{\B}{\mathbb B}
\newcommand{\C}{\mathbb C}
\newcommand{\D}{\mathbb D}
\newcommand{\R}{\mathbb R}
\newcommand{\V}{\mathbb V}
\newcommand{\N}{\mathbb N}
\newcommand{\M}{\mathbb M}
\newcommand{\K}{\mathbb K}
\newcommand{\F}{\mathbb F}
\newcommand{\X}{\mathbb X}
\newcommand{\Ha}{\mathbb H}
\newcommand{\Sa}{\mathbb S}
\newcommand{\mcA}{\mathcal A}
\newcommand{\mcN}{\mathcal N}
\newcommand{\mcV}{\mathcal V}
\newcommand{\mcH}{\mathcal H}
\newcommand{\mcB}{\mathcal B}
\newcommand{\mcC}{\mathcal C}
\newcommand{\mcD}{\mathcal D}
\newcommand{\mcE}{\mathcal E}
\newcommand{\mcF}{\mathcal F}
\newcommand{\mcP}{\mathcal P}
\newcommand{\mcQ}{\mathcal Q}
\newcommand{\mcL}{\mathcal L}
\newcommand{\mcJ}{\mathcal J}
\newcommand{\mcR}{\mathcal R}
\newcommand{\mcS}{\mathcal S}
\newcommand{\mcX}{\mathcal X}
\newcommand{\mcZ}{\mathcal Z}
\newcommand{\mcM}{\mathcal M}
\newcommand{\mfB}{\mathfrak B}
\newcommand{\mfF}{\mathfrak F}
\newcommand{\mfD}{\mathfrak D}
\newcommand{\mfL}{\mathfrak L}
\newcommand{\mfO}{\mathfrak O}
\newcommand{\mfS}{\mathfrak S}
\newcommand{\mfX}{\mathfrak X}
\newcommand{\mfR}{\mathfrak R}
\newcommand{\mfT}{\mathfrak T}
\newcommand{\mfM}{\mathfrak M}
\newcommand{\ov}{\overline}
\newcommand{\mb}{\makebox}
\newcommand{\es}{\emptyset}
\newcommand{\ci}{\subseteq}
\newcommand{\rec}[1]{\frac{1}{#1}}
\newcommand{\ld}{\ldots}
\newcommand{\ds}{\displaystyle}
\newcommand{\f}{\frac}
\newcommand{\tf}{\tfrac}
\newcommand{\et}[2]{#1_{1}, #1_{2}, \ldots, #1_{#2}}
\newcommand{\al}{\alpha}
\newcommand{\be}{\beta}
\newcommand{\G}{\Gamma}
\newcommand{\g}{\gamma}
\newcommand{\e}{\varepsilon}
\newcommand{\ph}{\phi}
\newcommand{\de}{\delta}
\newcommand{\De}{\Delta}
\newcommand{\la}{\lambda}
\newcommand{\La}{\Lambda}
\newcommand{\m}{\mu}
\newcommand{\om}{ \omega}
\newcommand{\Om}{\Omega}
\newcommand{\Si}{\Sigma}
\newcommand{\s}{\sigma}
\newcommand{\sgn}{\operatorname{sgn}}
\newcommand{\lb}{\label}
\newcommand{\rf}{\ref}
\newcounter{cnt1}
\newcounter{cnt2}
\newcounter{cnt3}
\newcommand{\blr}{\begin{list}{$($\roman{cnt1}$)$}
 {\usecounter{cnt1} \setlength{\topsep}{0pt}
 \setlength{\itemsep}{0pt}}}
\newcommand{\bla}{\begin{list}{$($\alph{cnt2}$)$}
 {\usecounter{cnt2} \setlength{\topsep}{0pt}
 \setlength{\itemsep}{0pt}}}
\newcommand{\bln}{\begin{list}{$($\arabic{cnt3}$)$}
 {\usecounter{cnt3} \setlength{\topsep}{0pt}
 \setlength{\itemsep}{0pt}}}
\newcommand{\el}{\end{list}}
\newtheorem{thm}{Theorem}[section]
\newtheorem{lem}[thm]{Lemma}
\newtheorem{cor}[thm]{Corollary}
\newtheorem{ex}[thm]{Example}
\newtheorem{Q}[thm]{Question}
\newtheorem{Def}[thm]{Definition}
\newtheorem{prop}[thm]{Proposition}
\newtheorem{rem}[thm]{Remark}
\newcommand{\Rem}{\begin{rem} \rm}
\newcommand{\bdfn}{\begin{Def} \rm}
\newcommand{\edfn}{\end{Def}}
\newcommand{\tx}{\text}
\newcommand{\TFAE}{the following are equivalent~: }
\newcommand{\etc}[3]{#1_{#3 1}, #1_{#3 2}, \ld, #1_{#3 #2}}
\newcommand{\ba}{\begin{array}}
\newcommand{\ea}{\end{array}}
\newcommand{\alto}[3]{\lt\{ \ba {ll}#1 & \mb{ if \quad}#2 \\ #3
 & \mb{ otherwise} \ea \rt.}
\newtheorem*{acknowledgements}{Acknowledgements}
\numberwithin{equation}{section}
\sloppy
\date{}
\begin{document}
\title{\bf Relativistic Newtonian Theory}
\author[Gill]{Tepper L. Gill}
\address[Tepper L. Gill]{ Department of EECS, Mathematics, and Computation Physics Laboratory, Howard University\\
Washington DC 20059 \\ USA, {\it E-mail~:} {\tt tgill@howard.edu}}

%thispagestyle{empty}
\subjclass{Primary (45) Secondary(46) }
\keywords{tensor product, relative tensor norm,  infinite-dimensional Laplacian}
\maketitle
\begin{abstract}  The objective of this paper is to derive a relativistic theory of Newtonian mechanics, which is compatible with Maxwell's theory by using the Einstein dual theory of special relativity.
\end{abstract}
\section{\bf{Dual Classical Theory}}
\subsection{Particle Clock}
To develop the dual classical theory, we assume a classical interacting particle defined on phase space with variables $( {\bf{x}}, {\bf{p}})$, clock $t$ and classical Hamiltonian $H=\sqrt{c^2{\bf{p}}^2+ m^2c^4}$, as seen by an observer $O$ in an inertial frame; and with variables $( {\bf{x}}', {\bf{p}}')$, clock $t'$ and classical Hamiltonian $H'=\sqrt{c^2{\bf{p}'}^2+ m^2c^4}$ as seen by an observer $O'$ in another inertial frame.       If ${\mathbf{w}}, {\mathbf{w}}'$ represent the particle velocities, let ${\gamma ^{ - 1}}\left( {\mathbf{w}}\right) = \sqrt {1 - {{{{\mathbf{w}}^2}} \mathord{\left/
 {\vphantom {{{{\mathbf{w}}^2}} {{c^2}}}} \right.
 \kern-\nulldelimiterspace} {{c^2}}}} $, ${\gamma ^{ - 1}}\left( {\mathbf{w}'}\right) = \sqrt {1 - {{{{\mathbf{w}'}^2}} \mathord{\left/
 {\vphantom {{{{\mathbf{w}'}^2}} {{c^2}}}} \right.
 \kern-\nulldelimiterspace} {{c^2}}}} $.  The classical proper time is defined by:
\beqn\lb{a1} 
d\tau  =\sqrt {1 - {{{{\mathbf{w}}^2}} \mathord{\left/
 {\vphantom {{{{\mathbf{w}}^2}} {{c^2}}}} \right.
 \kern-\nulldelimiterspace} {{c^2}}}}dt,\quad {\mathbf{w}} = \frac{{d{\mathbf{x}}}}{{dt}},\quad d{\tau^2} = d{t^2} - \tfrac{1}{{{c^2}}}d{{\mathbf{x}}^2}.
\eeqn  
\beqn\lb{b1}
d\tau  =\sqrt {1 - {{{{\mathbf{w}'}^2}} \mathord{\left/
 {\vphantom {{{{\mathbf{w}'}^2}} {{c^2}}}} \right.
 \kern-\nulldelimiterspace} {{c^2}}}}dt,\quad {\mathbf{w}'} = \frac{{d{\mathbf{x}'}}}{{dt'}},\quad d{\tau^2} = d{{t'}^2} - \tfrac{1}{{{c^2}}}d{{\mathbf{x}'}^2}.  
\eeqn
This leads to the Minkowski formulation of Einstein's special theory of relativity.

If we rearranging the last term in the above equations, we get:
\begin{equation}\lb{c}
\quad \quad d{t^2} = d{\tau^2} + \tfrac{1}{{{c^2}}}d{{\mathbf{x}}^2}, \Rightarrow cdt = \left( {\sqrt {{{\mathbf{u}}^2} + {c^2}} } \right)d\tau ,\quad {\mathbf{u}} = \frac{{d{\mathbf{x}}}}{{d\tau }}= {\gamma}({\bf{w}}){{\mathbf{w}}}.
\end{equation}
and
\begin{equation}\lb{d}
\quad \quad d{t'^2} = d{\tau^2} + \tfrac{1}{{{c^2}}}d{{\mathbf{x}'}^2}, \Rightarrow cdt' = \left( {\sqrt {{{\mathbf{u}'}^2} + {c^2}} } \right)d\tau ,\quad {\mathbf{u}'} = \frac{{d{\mathbf{x}'}}}{{d\tau }}= {\gamma}({\bf{w}'}){{\mathbf{w}'}}.
\end{equation}
If we let $b={\sqrt {{{\mathbf{u}}^2} + {c^2}} }$, $b'={\sqrt {{{\mathbf{u}'}^2} + {c^2}} }$; the second terms in the above equations become $cdt =bd\tau$ and $cdt' =b'd\tau$.  This leads to our first set of identities:
\begin{equation}\label{e}
\frac{1}{c}\frac{d}{{dt}} \equiv \frac{1}{{{b}}}\frac{d}{{d{\tau}}}, \quad \frac{1}{c}\frac{d}{{dt'}} \equiv \frac{1}{{{b'}}}\frac{d}{{d{\tau}}}
\end{equation}
These identities provide an  equally valid way to define the relationship between the (particle or) proper time and the observer time for particle dynamics. 

If we apply them to ${\bf{x}}$ and ${\bf{x}'}$, we obtain our second new identities, showing that the transformation leaves the tangent space invariant:
\begin{equation}\label{f}
\frac{{{{\mathbf{w}}}}}{c} = \frac{1}{c}\frac{{d{{\mathbf{x}}}}}{{dt}} \equiv \frac{1}{{{b}}}\frac{{d{{\mathbf{x}}}}}{{d\tau }} = \frac{{{{\mathbf{u}}}}}{{{b}}}, \quad \frac{{{{\mathbf{w}'}}}}{c} = \frac{1}{c}\frac{{d{{\mathbf{x}'}}}}{{dt'}} \equiv \frac{1}{{{b'}}}\frac{{d{{\mathbf{x}'}}}}{{d\tau }} = \frac{{{{\mathbf{u}'}}}}{{{b'}}}. 
\end{equation}

  
\begin{rem} The mapping  $t \to \tau$ (respectively $t' \to \tau$) is a member of the family of contact groups, often used in celestial mechanics (see \cite{9}). Contact transformations are sometimes called tangency transformations in mechanics, because they leave invariant the tangent at the point of contact.    
\end{rem}

\begin{thm}The phase space variables: (${\bf{x}}, {\bf{p}}$) are also left invariant by the contact mapping  $t \to \tau$ (respectively (${\bf{x}'}, {\bf{p}'}$) by the contact mapping $t' \to \tau$). 
\end{thm}
\begin{proof} Clearly the position vectors are invariant, so we need only check the momentum variables.  If $m$ is the rest mass, it is sufficient to notice that ${\mathbf{p}} = m\gamma \left( {\mathbf{w}} \right){\mathbf{w}}$ and ${\mathbf{u}} = \gamma \left( {\mathbf{w}} \right){\mathbf{w}}$, so that ${\mathbf{p}} = m{\mathbf{u}}$ and, by the same observation ${\mathbf{p}'} = m{\mathbf{u}'}$.
\end{proof}
The new particle coordinates are $({\bf{x}}, \tau)$ (respectively $({\bf{x}'}, \tau)$).  In this representation, the position is uniquely defined by the observers, while $\tau$ is uniquely defined by the particle. 
\subsection{\bf Dual Particle Theory}
In the conventional formulation of quantum theory, the Hamiltonian $H$ is the generator of observer time translations (i.e, $t$ translations). We now seek to identify the Hamiltonian $K$ which will generate  proper-time or $\tau$ translations. To see how this may be done, let $W$ be any classical phase space observable so that the Poisson bracket defines Hamilton's equations in the $O$ frame by: (recall $H = \sqrt {{c^2}{{\mathbf{p}}^2} + {m^2}{c^4}} $)
\beqn
\frac{{dW}}{{dt}} = \left\{ {H,W} \right\} = \frac{{\partial H}}{{\partial {\mathbf{p}}}}\frac{{\partial W}}{{\partial {\mathbf{x}}}} - \frac{{\partial H}}{{\partial {\mathbf{x}}}}\frac{{\partial W}}{{\partial {\mathbf{p}}}}.
\eeqn
To represent the dynamics via the proper time of the particle, we first recall that, for any particle of rest mass $m$, the Hamiltonian  can be represented as $H = mc^2\g({\bf{w}})$, so that $\g({\bf{w}}) = H/mc^2$. This means that we can also represent the proper time by $d\tau  = ({{mc^2 } \mathord{\left/ {\vphantom {{mc^2 } H}} \right. \kern-\nulldelimiterspace} H})dt$.  Using the chain rule, we have
\[
\frac{{dW}}{{d\tau }} = \frac{{dt}}{{d\tau }}\frac{{dW}}{{dt}} = \frac{H}
{{mc^2 }}\left\{ {H,W} \right\}.
\]

Using the invariant rest energy $mc^2$, we determine the canonical proper-time Hamiltonian $K$ such that:
\[
\left\{ {K,W} \right\} = \frac{H}{{mc^2 }}\left\{ {H,W} \right\},\quad \left. K \right|_{{\mathbf{p}} = 0}  = \left. H \right|_{{\mathbf{p}} = 0}  = mc^2. 
\]
From 
\[
\begin{gathered} 
  \left\{ {K,W} \right\} = \left[ {\frac{H}{{mc^2 }}\frac{{\partial H}}{{\partial {\mathbf{p}}}}} \right]\frac{{\partial W}}{{\partial {\mathbf{x}}}} - \left[ {\frac{H}
{{mc^2 }}\frac{{\partial H}}{{\partial {\mathbf{x}}}}} \right]\frac{{\partial W}}
{{\partial {\mathbf{p}}}} \hfill \\
  {\text{          }} = \frac{\partial }{{\partial {\mathbf{p}}}}\left[ {\frac{{H^2 }}
{{2mc^2 }} + a} \right]\frac{{\partial W}}{{\partial {\mathbf{x}}}} - \frac{\partial }
{{\partial {\mathbf{x}}}}\left[ {\frac{{H^2 }}{{2mc^2 }} + a'} \right]\frac{{\partial W}}
{{\partial {\mathbf{p}}}}, \hfill \\ 
\end{gathered} 
\]
we see that $a = a' = \tfrac{1}{2}mc^2$.  Thus, assuming no explicit time dependence, we have:
\beqn{\lb{2.4}}
K = \frac{{H^2 }}{{2mc^2 }} + \frac{{mc^2 }}{2},\quad {\text{and }}\quad \frac{{dW}}
{{d\tau }} = \left\{ {K,W} \right\}.
\eeqn 
\subsection{Interaction}
Since $\tau$ remains invariant during interaction, we assume  $K$ also remains invariant. For the classical relativistic case, there are two possible Hamiltonians:
\[
\begin{gathered}
  {H_1} = \sqrt {{c^2}{{\bs{\pi}}^2} + {{\left( {m{c^2} + V} \right)}^2}}  \hfill \\
  {H_2} = \sqrt {{c^2}{{\bs{\pi}}^2} + {m^2}{c^4}}  + V \hfill \\ 
\end{gathered} 
\]
For the Newtonian-Maxwell relativistic unification, we use the first one. If $\sqrt {{c^2}{{\bf{p}}^2}  + m^2 c^4 }  \to \sqrt {c^2 {\bs{\pi}}^2  + (m c^2 + V)^2}  $,  $K$ becomes:
\beqn
K = \frac{{{\bs\pi ^2}}}{{2m}} + m{c^2} + V + \frac{{{V^2}}}{{2m{c^2}}}.
\eeqn
We consider the simple case of a particle moving under the action of a central force, with potential 
$V=- \tf{\la}{r},\; \la>0$.  From $\tfrac{{\partial K}}{{\partial {\mathbf{p}}}} = {\mathbf{u}}$, we have  ${\mathbf{u}} = \tfrac{{\mathbf{p}}}{{m}}$ and ${\mathbf{a}} = \tfrac{{ \mathbf{\dot p}}}{{m}}$.
Solving for ${\bf{p}}, \;{\bf{p}}^2$ and $\dot{\bf{p}}$, we have 
${\mathbf{p}} = m{\mathbf{u}}, \; {{\mathbf{p}}^2} = {m^2}{{\mathbf{u}}^2}$, and ${\mathbf{\dot p}} = m{\mathbf{a}}$.    
Finally,
from $- \tfrac{{\partial K}}{{\partial {\mathbf{x}}}} = {\mathbf{\dot p}}$, we have
\beqn\lb{h3}
{\mathbf{\dot p}} =m{\mathbf{a}}= - \nabla V-\nabla V \tfrac{V}{{m{c^2}}}\; \Rightarrow \; {\mathbf{a}}=- \tf{\nabla V}{m}(1+ \tfrac{V}{{m{c^2}}}).
\eeqn
In order to see the physical meaning of this equation, we  note that since $V$ is attractive, $-\bs\nabla V (V/mc^2)>0$, so there is a value $r_0$ for which  $-\bs\nabla V   -\bs\nabla V (V/mc^2)=0$.  Thus, this equation is attractive for $r>r_0$, is zero at $r=r_0$ and is repulsive when $0< r< r_0$. (For example, if $V$ is the potential between an electron and proton, $r=r_0$ is the classical electron radius.)  We interpret  $0< r< r_0$ as a fixed region of repulsion, so that the singularity $r=0$ is impossible to reach at the classical (relativistic Newtonian) level.    This means that the classical principle of impenetrability, namely that no two particles can occupy the same space at the same time occurs naturally in the relativistic Newtonian theory.
\begin{rem} We expect that this effect should have some impact on the galactic and cosmological scale. 
\end{rem}
\section{The Sun-Mercury System}
In this section assume an external observer views the Sun-Mercury system.  In this case, we have a two body problem, the line of force is the same but the equations of motion are:
\beqn\lb{h4}
\begin{gathered}
  m{{\mathbf{a}}_m} =  - \tfrac{{GMm}}{{{r^2}}}\left( {1 - \tfrac{{GM}}{c^2r}} \right){{\mathbf{e}}_r} \Rightarrow \quad {{\mathbf{a}}_m} =  - \tfrac{{GM}}{{{r^2}}}\left( {1 - \tfrac{{GM}}{c^2r}} \right){{\mathbf{e}}_r} \hfill \\
  M{{\mathbf{a}}_s} = \tfrac{{GMm}}{{{r^2}}}\left( {1 - \tfrac{{Gm}}{c^2r}} \right){{\mathbf{e}}_r} \Rightarrow \quad {{\mathbf{a}}_s} = \tfrac{{Gm}}{{{r^2}}}\left( {1 - \tfrac{{Gm}}{c^2r}} \right){{\mathbf{e}}_r}. \hfill \\ 
\end{gathered} 
\eeqn
We can now find the relative acceleration $\bf{a}$ from:
\beqn
\begin{gathered}
  {\mathbf{a}} = {{\mathbf{a}}_m} - {{\mathbf{a}}_s} =  - \frac{{GM}}{{{r^2}}}\left( {1 - \frac{{GM}}{{{c^2}r}}} \right){{\mathbf{e}}_r} - \frac{{Gm}}{{{r^2}}}\left( {1 - \frac{{Gm}}{{{c^2}r}}} \right){{\mathbf{e}}_r} \hfill \\
   =  - \frac{{G\left( {M + m} \right)}}{{{r^2}}}\left[ {1 - \frac{{{G}\left( {{M^2} + {m^2}} \right)}}{{\left( {M + m} \right){c^2}r}}} \right]{{\mathbf{e}}_r}. \hfill \\ 
\end{gathered} 
\eeqn
Thus, the effective gravitational force on the planet Mercury is 
\[{{{F}}_e} =  - \frac{{G\left( {M + m} \right)m}}{{{r^2}}}\left[ {1 - \frac{{{G}\left( {{M^2} + {m^2}} \right)}}{{\left( {M + m} \right){c^2}r}}} \right].\]
If $r_0$ is the orbital radius, we can determine the orbital period by equating the effective gravitational force to the centripetal force and solving for the dual velocity of rotation $u_d$: 
 \beqn\lb{h5}
\begin{gathered}
  \frac{{G\left( {M + m} \right)m}}{{r_0^2}}\left[ {1 - \frac{{{G}\left( {{M^2} + {m^2}} \right)}}{{\left( {M + m} \right){c^2}{r_0}}}} \right] = \frac{{mu_d^2}}{{{r_0}}} \hfill \\
   \Rightarrow {u_d} = \frac{{{{\left[ {G\left( {M + m} \right)} \right]}^{\tfrac{1}{2}}}}}{{r_0^{\tfrac{1}{2}}}}{\left[ {1 - \frac{{{G}\left( {{M^2} + {m^2}} \right)}}{{\left( {M + m} \right){c^2}{r_0}}}} \right]^{\tfrac{1}{2}}}, \hfill \\ 
\end{gathered}
\eeqn
 The (non-relativistic) Newtonian period is defined by 
 \beqn\lb{h6}
T_0=\f{2\pi r_0}{u_0}= \f{2\pi r_0^{\tf{3}{2}}}{(GM)^{\tf{1}{2}}} 
\eeqn
 with corresponding angular velocity $\om_0=2\pi/T_0$. Thus,   the angular distance that Mercury sweeps during the Newtonian orbital period $T_0$ in radians per revolution is
$\varphi_0 = \om_0 T_0=2\pi$.
From equation (2.3), we see that  the relativistic or dual Newtonian period is given by
 \beqn\lb{h7}
\begin{gathered}
  {T_d} = \frac{{2\pi }}{{{u_d}}} = \frac{{2\pi r_o^{\tfrac{3}{2}}}}{{{{\left[ {GM} \right]}^{\tfrac{1}{2}}}}}{\left[ {\left( {1 + \tfrac{m}{M}} \right)\left( {1 - \tfrac{{G\left( {{M^2} + {m^2}} \right)}}{{\left( {M + m} \right){c^2}r_0}}} \right)} \right]^{ - \tfrac{1}{2}}} \hfill \\
   = {T_0}{\left[ {\left( {1 + \tfrac{m}{M}} \right)\left( {1 - \tfrac{{G\left( {1 + \tfrac{{{m^2}}}{{{M^2}}}} \right)}}{{\left( {1 + \tfrac{m}{M}} \right){c^2}r_0}}} \right)} \right]^{ - \tfrac{1}{2}}}. \hfill \\ 
\end{gathered}
\eeqn
The dual angular velocity $\omega _d$ is
\beqn\lb{h8}
\begin{gathered}
  {\omega _d} = \frac{{2\pi }}{{{T_d}}} = \frac{{2\pi }}{{{T_0}}}{\left[ {\left( {1 + \tfrac{m}{M}} \right)\left( {1 - \tfrac{{G\left( {{M^2} + {m^2}} \right)}}{{\left( {M + m} \right){c^2}r_0}}} \right)} \right]^{\tfrac{1}{2}}}  \hfill \\
  = {\omega _0}{\left[ {\left( {1 + \tfrac{m}{M}} \right)\left( {1 - \tfrac{{G\left( {{M^2} + {m^2}} \right)}}{{\left( {M + m} \right){c^2}r_0}}} \right)} \right]^{\tfrac{1}{2}}} \hfill \\
= {\omega _0}\left[ {{{\left( {1 + \tfrac{m}{M}} \right)}^{\tfrac{1}{2}}}{{\left( {1 - \tfrac{{G\left( {{M^2} + {m^2}} \right)}}{{\left( {M + m} \right){c^2}r_0}}} \right)}^{\tfrac{1}{2}}}} \right]  \hfill \\
\end{gathered}
\eeqn
The second term in brackets can be reduced to obtain:
\[\begin{gathered}
  1 - \frac{{G\left( {{M^2} + {m^2}} \right)}}{{\left( {M + m} \right){c^2}{r_0}}} = 1 - \frac{{G{{\left( {M + m} \right)}^2} - 2GMm}}{{\left( {M + m} \right){c^2}{r_0}}} \hfill \\
   = 1 + \frac{{2GMm}}{{\left( {M + m} \right){c^2}{r_0}}} - \frac{{G\left( {M + m} \right)}}{{{c^2}{r_0}}} \hfill \\
   = 1 - \frac{{GM\left( {1 + \tfrac{m}{M}} \right)}}{{{c^2}{r_0}}} + \frac{{2G\mu }}{{{c^2}{r_0}}}\approx 1 - \frac{{GM\left( {1 + \tfrac{m}{M}} \right)}}{{{c^2}{r_0}}}, \hfill \\ 
\end{gathered} 
\]
where $\mu  = Mm/(M + m)$ is the reduced mass.

We now have three possible angles that Mercury can sweep during the period $T_0$, based on the following:
\begin{enumerate}
\item Follow Corda, using $\omega _c=\omega _0{\left(1 + \tfrac{m}{{M}} \right)^{\tf{1}{2}}} \approx \omega _0{\left(1 + \tfrac{m}{{2M}} \right)}$, so that ${\varphi _c} = 2\pi \left( {1 + \tfrac{m}{{2M}}} \right)$.  Thus, as ${\varphi _0}=2\pi$ we obtain:
\[
\Delta {\varphi _c} = {\varphi _c} - 2\pi  = 2\pi \left( {1 + \tfrac{m}{{2M}}} \right) - 2\pi  = \frac{{\pi m}}{M},
\]
leading to a value of $44.39"$ per tropical century.
\item Improve on Corda's result by using the relativistic Newtonian result:
\beqn
\omega _{d}={\omega _0}{\left[ {\left( {1 + \tfrac{m}{M}} \right)\left( {1 - \tfrac{{G\left( {{M^2} + {m^2}} \right)}}{{\left( {M + m} \right){c^2}r_0}}} \right)} \right]^{\tfrac{1}{2}}}.
\eeqn
Thus, 
\beqn
{\varphi _d} = 2\pi {\left[ {\left( {1 + \tfrac{m}{M}} \right)\left( {1 - \frac{{G\left( {{M^2} + {m^2}} \right)}}{{\left( {M + m} \right){c^2}{r_0}}}} \right)} \right]^{\tfrac{1}{2}}}
\eeqn
and $\Delta {\varphi _d} = {\varphi _d} - 2\pi$.
\item Improve by using the relativisitic approximation:
\[
{\omega _{{d_1}}} \simeq {\omega _0}\left( {1 + \tfrac{m}{{2M}}} \right)\left[ {1 - \frac{{GM\left( {1 + \tfrac{m}{M}} \right)}}{{2{c^2}{r_0}}}} \right] \simeq {\omega _0}\left[ {1 + \frac{m}{{2M}} - \frac{{GM\left( {1 + \tfrac{m}{M}} \right)}}{{2{c^2}{r_0}}}} \right],
\]
so that we obtain:
\[
 \Delta {\varphi _{{d_1}}} = {\varphi _{_{{d_1}}}} - 2\pi  = 2\pi \left[ {1 + \frac{m}{{2M}} - \frac{{GM\left( {1 + \tfrac{m}{M}} \right)}}{{2{c^2}{r_0}}}} \right] - 2\pi .
\]
\end{enumerate}
The second and third possibilities are of main interest as each gives a smaller value than  $\omega _c$ (due to Corda).  To see how this is possible, the unperturbed angle Mercury sweeps is ${\varphi _0}=2\pi$ and the Corda improvement is:
\beqn
\Delta {\varphi _c} = {\varphi _c} - 2\pi  = \tfrac{{\pi m}}{M}\left( {1 - \tfrac{{m}}{{4M}}} \right) \; {\rm{implies}}\; \Delta {\varphi _c} = 44.39".
\eeqn
This well approximates the value per tropical century of $42.98"$  given by general relativity and the well known observational value of $43"$. 

\begin{thebibliography}{99}
\small
%\addcontentsline{toc}{chapter}{\protect\numberline{}
\bibliographystyle{amsalpha}

\bibitem[1]{1} T. L. Gill and W. W. Zachary, {\it Foundations for relativistic quantum theory I: Feynman's operator calculus and the Dyson conjectures}, Journal of  Mathematical Physics  {\bf 43} (2002), 69-93.
\bibitem[2]{2} T.L. Gill, and G. Ares de Parga,{\it The Einstein Dual Theory of Relativity},  Advanced Studies in Theoretical Physics {\bf 13}(8) (2019) 337-377.
\bibitem[3]{3} P. A. M. Dirac, Proc. Roy. Soc (London) {\bf{A117}} (1928) 610, {\bf{A118}} (1928) 351.
\bibitem[4]{4} J.A. Wheeler and R.P. Feynman, Interaction with the absorber as the mechanism of radiation, {\it{ Rev. Mod. Phys.}}, {\bf{17}} (1949), 157-181.
\bibitem[5]{5} T.L. Gill, W.W. Zachary, and J. Lindesay,  {\it The Classical Electron Problem},   Foundations of Phys. {\bf{31}} (2001) 1299-1354.
\bibitem[6]{6} T. L. Gill and W. W. Zachary, {\it Two Mathematically Equivalent Versions of Maxwell's Equations},   Foundations of Phys. {\bf{41}} (2011) 99-128. 
\bibitem[7]{7}T.L. Gill, T. Morris and S. K. Kurtz  Universal Journal of Physics and Application {\bf 3}(1)(2015) 24-40.
\bibitem[8]{8} R. M. Santilli, {\it Isonumbers and genonumbers of dimension 1, 2, 4, 8, their isoduals, and pseudoduals, and ``hidden numbers" of dimension 3, 5, 6, 7}, Algebras, Groups and Geometries, Vol. {\bf 10}, 273-322, 1993.
\bibitem[9]{9}E. T. Whittaker,  {\it A Treatise on the Analytical Dynamics of Particles and rigid bodies}, Cambridge, University Press, Cambridge, 1917.
\bibitem[10]{10} T. L. Gill and J. Lindesay, {\it Canonical Proper Time Formulation of Relativistic Particle Dynamics}, Int. J. of Theor. Phys.  {\bf 32} (1993), 2087-2098.
\bibitem[11]{11} T. L. Gill, {\it The Square-Root Operator,  Proper-Time and a Particle Interpretation for the Klein-Gordon Equation}, FERMILAB-Pub-82/60-THY (1982).
\bibitem[12]{12} T.L. Gill, W.W. Zachary,  {\it Analytic representation of the square-root operator},   J. Phys. A: Math. Gen. {\bf 38} (2005) 1-18.
\bibitem[13]{13} T.L. Gill, W.W. Zachary, and Marcus Alfred,  {\it Analytic representation of the Dirac equation},   J. Phys. A: Math. Gen. {\bf 38} (2005) 1-22.
\bibitem[14]{14} S. S. Schweber, {\it QED and the Men Who Made It: Dyson, Feynman, Schwinger and Tomonaga }, Princeton University Press, Princeton NJ, (1994).
\end{thebibliography}
\end{document} 
